\chapter{Meet the Bot We Will Build}

\begin{goalbox}
Understand the finished result, the few tools involved, and the project layout.
\end{goalbox}

Our goal is intentionally small: a Telegram bot that is easy to finish and easy
to understand. It will:

\begin{itemize}
\item welcome a user who sends \texttt{/start};
\item explain itself when the user sends \texttt{/help};
\item reply to ordinary text with \texttt{You said:} followed by that text;
\item ask for text when it receives a photo, voice note, or similar content.
\end{itemize}

For example, when Alireza sends \texttt{My first bot}, the conversation is:

\begin{resultbox}
\textbf{Alireza:} My first bot\\
\textbf{Bot:} You said: My first bot
\end{resultbox}

\section{The four pieces}

Telegram shows the conversation. BotFather creates the bot account and gives its
secret token. Python runs our instructions. The \texttt{pyTelegramBotAPI}
package provides the \texttt{telebot} module that connects those instructions to
Telegram. This tutorial uses polling: while the program runs, it regularly asks
Telegram whether a new message has arrived.

\section{A small, readable layout}

\begin{lstlisting}[language={},numbers=none]
telbot-simple-bot/
|-- code/
|   |-- bot.py
|   |-- requirements.txt
|   |-- test_bot.py
|   `-- README.md
|-- documentation/
|   |-- English/
|   `-- Persian/
|-- LINKEDIN_POST.md
`-- README.md
\end{lstlisting}

Only \texttt{code/bot.py} is needed to run the bot after the dependency is
installed. The other files teach, test, or document it.

\begin{checkpointbox}
You know exactly what the finished bot will do. No token, package, or account has
been changed yet.
\end{checkpointbox}
